% GENERATED FILE. Edit canonical lessons, then run npm run generate:books.
% canonical-source-hash: fnv1a64:c544c3fb21956349
% canonical-lessons: FR-C271-peu-a-peu, FR-C271-progressivement, FR-C271-au-fur-et-a-mesure, FR-C271-a-la-longue, FR-C271-tot-ou-tard

\chapter{Peu à peu}
\label{ch:fr-peu-a-peu}
\hlchaptermodality{271}

\begin{chapteropening}
\textbf{By the end of this chapter:} \emph{I can say that something changed little by little, gradually, as things went along, in the long run, or sooner or later.}

Complete the last lesson of chapter 271: I can say that something changed little by little, gradually, as things went along, in the long run, or sooner or later.
\end{chapteropening}

\section[peu à peu]{peu à peu — little by little}

\label{lesson:FR-C271-peu-a-peu}



\begin{quote}
Before the new one: say the French for at the time, then the French for straight after.
\end{quote}

\subsection*{You'll want to know: peu à peu}
\textbf{peu à peu} — \textquotedblleft{}little by little\textquotedblright{}.

\begin{quote}
Peu à peu, la ville s'est réveillée.
\end{quote}

\begin{quote}
\emph{Little by little, the town woke up.}
\end{quote}

The same as \emph{petit à petit}.

\begin{cousinweb}
\emph{peu} is inherited from Latin \emph{paucum} \textquotedblleft{}a little\textquotedblright{}, the neuter of \emph{paucus} \textquotedblleft{}few\textquotedblright{} (Lewis \& Short).
\end{cousinweb}

\begin{grammarlens}[title={What you've built}]
One more word for telling, explaining and describing at B1.
\end{grammarlens}

\subsection*{Your turn}
\begin{itemize}
\raggedright
  \item \emph{Say it:} \emph{peu à peu}
  \item \emph{Say it:} \emph{peu à peu}, once more
  \item \emph{Say it:} a sentence of your own with \emph{peu à peu}
  \item \emph{Recall:} say the French for at the time, then the French for straight after, then say \emph{peu à peu} again
\end{itemize}

\begin{tcolorbox}[breakable,colback=teal!4,colframe=teal!35!black,title={Before you move on}]
What does peu à peu mean? (\textquotedblleft{}Little by little\textquotedblright{}.) Say it once more.
\end{tcolorbox}
\section[progressivement]{progressivement — gradually, progressively}

\label{lesson:FR-C271-progressivement}



\begin{quote}
Before the new one: say the French for straight after, then the French for little by little.
\end{quote}

\subsection*{You'll want to know: progressivement}
\textbf{progressivement} — \textquotedblleft{}gradually, progressively\textquotedblright{}.

\begin{quote}
Il fait progressivement plus chaud.
\end{quote}

\begin{quote}
\emph{It is gradually getting warmer.}
\end{quote}

Built on \emph{progressif}; a neutral word, common in reports.

\begin{cousinweb}
\emph{progressivement} is built on \emph{progressif}, a learned borrowing from Latin \emph{prōgressīvus} (EtymDB), built on \emph{prōgressus} \textquotedblleft{}a going forward\textquotedblright{}, from \emph{prōgredī} \textquotedblleft{}to step forward\textquotedblright{} (Lewis \& Short).
\end{cousinweb}

\begin{grammarlens}[title={What you've built}]
One more word for telling, explaining and describing at B1.
\end{grammarlens}

\subsection*{Your turn}
\begin{itemize}
\raggedright
  \item \emph{Say it:} \emph{progressivement}
  \item \emph{Say it:} \emph{progressivement}, once more
  \item \emph{Say it:} a sentence of your own with \emph{progressivement}
  \item \emph{Recall:} say the French for straight after, then the French for little by little, then say \emph{progressivement} again
\end{itemize}

\begin{tcolorbox}[breakable,colback=teal!4,colframe=teal!35!black,title={Before you move on}]
What does progressivement mean? (\textquotedblleft{}Gradually, progressively\textquotedblright{}.) Say it once more.
\end{tcolorbox}
\section[au fur et à mesure]{au fur et à mesure — as one goes along, bit by bit}

\label{lesson:FR-C271-au-fur-et-a-mesure}



\begin{quote}
Before the new one: say the French for little by little, then the French for gradually.
\end{quote}

\subsection*{You'll want to know: au fur et à mesure}
\textbf{au fur et à mesure} — \textquotedblleft{}as one goes along, bit by bit\textquotedblright{}.

\begin{quote}
Il apprend la langue au fur et à mesure.
\end{quote}

\begin{quote}
\emph{He learns the language as he goes along.}
\end{quote}

With a verb: \emph{au fur et à mesure que}: \emph{au fur et à mesure qu'on avance}.

\begin{cousinweb}
\emph{mesure} is inherited from Latin \emph{mēnsūra} \textquotedblleft{}a measuring\textquotedblright{} (EtymDB). \emph{Fur} is an old word for \textquotedblleft{}rate, price\textquotedblright{}, which EtymDB traces to Latin \emph{forum} \textquotedblleft{}market place\textquotedblright{}.
\end{cousinweb}

\begin{grammarlens}[title={What you've built}]
One more word for telling, explaining and describing at B1.
\end{grammarlens}

\subsection*{Your turn}
\begin{itemize}
\raggedright
  \item \emph{Say it:} \emph{au fur et à mesure}
  \item \emph{Say it:} \emph{au fur et à mesure}, once more
  \item \emph{Say it:} a sentence of your own with \emph{au fur et à mesure}
  \item \emph{Recall:} say the French for little by little, then the French for gradually, then say \emph{au fur et à mesure} again
\end{itemize}

\begin{tcolorbox}[breakable,colback=teal!4,colframe=teal!35!black,title={Before you move on}]
What does au fur et à mesure mean? (\textquotedblleft{}As one goes along, bit by bit\textquotedblright{}.) Say it once more.
\end{tcolorbox}
\section[à la longue]{à la longue — in the long run, eventually}

\label{lesson:FR-C271-a-la-longue}



\begin{quote}
Before the new one: say the French for gradually, then the French for as one goes along.
\end{quote}

\subsection*{You'll want to know: à la longue}
\textbf{à la longue} — \textquotedblleft{}in the long run, eventually\textquotedblright{}.

\begin{quote}
À la longue, on s'habitue au bruit.
\end{quote}

\begin{quote}
\emph{In the long run, you get used to the noise.}
\end{quote}

Here \emph{la longue} is a noun: a long stretch of time.

\begin{cousinweb}
\emph{longue} is the feminine of \emph{long}, inherited from Latin \emph{longus} \textquotedblleft{}long\textquotedblright{} (Lewis \& Short).
\end{cousinweb}

\begin{grammarlens}[title={What you've built}]
One more word for telling, explaining and describing at B1.
\end{grammarlens}

\subsection*{Your turn}
\begin{itemize}
\raggedright
  \item \emph{Say it:} \emph{à la longue}
  \item \emph{Say it:} \emph{à la longue}, once more
  \item \emph{Say it:} a sentence of your own with \emph{à la longue}
  \item \emph{Recall:} say the French for gradually, then the French for as one goes along, then say \emph{à la longue} again
\end{itemize}

\begin{tcolorbox}[breakable,colback=teal!4,colframe=teal!35!black,title={Before you move on}]
What does à la longue mean? (\textquotedblleft{}In the long run, eventually\textquotedblright{}.) Say it once more.
\end{tcolorbox}
\section[tôt ou tard]{tôt ou tard — sooner or later}

\label{lesson:FR-C271-tot-ou-tard}



\begin{quote}
Before the new one: say the French for as one goes along, then the French for in the long run.
\end{quote}

\subsection*{You'll want to know: tôt ou tard}
\textbf{tôt ou tard} — \textquotedblleft{}sooner or later\textquotedblright{}.

\begin{quote}
Tôt ou tard, il va comprendre.
\end{quote}

\begin{quote}
\emph{Sooner or later, he is going to understand.}
\end{quote}

The order is fixed: \emph{tôt} first.

\begin{cousinweb}
\emph{tard} is inherited from Latin \emph{tardē} \textquotedblleft{}slowly, late\textquotedblright{} (Lewis \& Short), from \emph{tardus} \textquotedblleft{}slow\textquotedblright{} (EtymDB); \emph{tôt} is usually traced to Latin \emph{tostum} \textquotedblleft{}scorched\textquotedblright{}.
\end{cousinweb}

\begin{grammarlens}[title={What you've built}]
One more word for telling, explaining and describing at B1.
\end{grammarlens}

\subsection*{Your turn}
\begin{itemize}
\raggedright
  \item \emph{Say it:} \emph{tôt ou tard}
  \item \emph{Say it:} \emph{tôt ou tard}, once more
  \item \emph{Say it:} a sentence of your own with \emph{tôt ou tard}
  \item \emph{Recall:} say the French for as one goes along, then the French for in the long run, then say \emph{tôt ou tard} again
\end{itemize}

\begin{tcolorbox}[breakable,colback=teal!4,colframe=teal!35!black,title={Before you move on}]
What does tôt ou tard mean? (\textquotedblleft{}Sooner or later\textquotedblright{}.) Say it once more.
\end{tcolorbox}
